% paper3_results.tex -- Results of Paper 3, input by papers/paper3_intercropping.tex.
% Every number is quoted from results/values.tex with \res{...}.
\section{Results}

In every block, each crop yielded less in the intercrop than as a sole crop,
as expected when each species occupies only half of the rows. The combined
production of the intercrop, however, exceeded that of the sole crops on an
equal-area basis at all three sites (Table~\ref{tab:p3ler}).

\input{\subfix{../tables/p3_ler}}

The mean total land equivalent ratio was \res{p3-ler-S1} at S1,
\res{p3-ler-S2} at S2 and \res{p3-ler-S3} at S3: to match one hectare of
intercrop, sole crops would have needed between \res{p3-ler-S1} and
\res{p3-ler-S3} hectares of land. Maize contributed the larger share at
every site, with partial ratios of \res{p3-lerm-S1}, \res{p3-lerm-S2} and
\res{p3-lerm-S3}; the partial ratios of cowpea were \res{p3-lerc-S1},
\res{p3-lerc-S2} and \res{p3-lerc-S3} (Figure~\ref{fig:p3ler}). The highest
ratio was found at the humid tropical site, where maize retained the largest
fraction of its sole-crop yield in the intercrop.

\input{\subfix{../figures/fig3_ler}}
